\documentclass[11pt]{amsart}
\usepackage{amsmath,amssymb,amsthm}
\usepackage[margin=1.1in]{geometry}
\usepackage{booktabs}
\usepackage[colorlinks=true,linkcolor=blue,citecolor=blue,urlcolor=blue]{hyperref}

\newtheorem{theorem}{Theorem}[section]
\newtheorem{lemma}[theorem]{Lemma}
\newtheorem{proposition}[theorem]{Proposition}
\newtheorem{corollary}[theorem]{Corollary}
\newtheorem{conjecture}[theorem]{Conjecture}
\newtheorem{hypothesis}[theorem]{Hypothesis}
\theoremstyle{definition}
\newtheorem{definition}[theorem]{Definition}
\newtheorem{remark}[theorem]{Remark}
\newtheorem{compfact}[theorem]{Computational Fact}

\DeclareMathOperator{\ord}{ord}
\DeclareMathOperator{\Geom}{Geom}
\newcommand{\N}{\mathbb{N}}
\newcommand{\Z}{\mathbb{Z}}
\newcommand{\R}{\mathbb{R}}
\newcommand{\E}{\mathbb{E}}
\newcommand{\Prob}{\mathbb{P}}
\newcommand{\e}{\mathrm{e}}
\newcommand{\ceil}[1]{\lceil #1\rceil}
\newcommand{\floor}[1]{\lfloor #1\rfloor}

\title{The 3-adic mixing rate of the Syracuse random variable:\\
reductions, a boundary-layer mechanism, and a sharp-rate conjecture}

\author{Hikaru Wakaura}
\address{QuantScape Inc., 4-11-18, Manshon-Shimizudai, Meguro, Tokyo, 153-0064, Japan}
\email{hikaruwakaura@gmail.com}
\date{\today}
\subjclass[2020]{Primary 11B83; Secondary 11K16, 37P10, 60H25}
\keywords{$3x+1$ problem, Syracuse random variable, Fourier decay,
$3$-adic equidistribution, Kesten chain, Furstenberg $\times 2\times 3$
rigidity, experimental mathematics}

\begin{document}

\begin{abstract}
Let $\mathbf{Syrac}$ be the Syracuse random variable, the $3$-adic limit of
$\sum_{j\ge 1} 3^{j-1}2^{-A_j}$ with $A_j$ the partial sums of i.i.d.\
$\Geom(1/2)$ exponents, and let
$\delta_m=\max_{3\nmid\xi}\,\bigl|\E\,\e(\xi\,\mathbf{Syrac}/3^m)\bigr|$ be
its maximal nontrivial Fourier coefficient --- the quantity whose decay
drives Tao's almost-all result for the Collatz map and any effective
version thereof. Tao proved $\delta_m\ll_A m^{-A}$ and explicitly flagged
the exponential improvement as possible but unpursued; no later work
supplies it. We prove: (1) since $2$ is a primitive root mod $3^m$, the
Fourier problem is exactly a one-dimensional lattice recursion over the
exponent of the frequency, with $\delta_m$ non-increasing and exactly
stationary after $m$ steps; (2) a first-passage renormalization identity
with the certified bound $|F(d,r)|\le 2^{-r}\delta_{d-1}$; (3) a weighted
contraction lemma whose optimal rate is exactly
$2^{-I_0}$, $I_0=(1-H(\theta))/\theta=0.0793186\ldots$, $\theta=\log_3 2$,
reducing the upper-bound problem to the control of frequencies near and
above the critical line $2^k\approx 3^m$; (4) the real Syracuse variable is
a Kesten chain with tail exponent exactly $1$; (5) an exact $2$-adic law:
over full multiplicative periods, the $2$-adic valuations of the signed
residues of $2^K \bmod 3^m$ inside a symmetric ball are geometrically
distributed, by an elementary two-line identity. Assembling (1)--(5) with
an audited weighted-$\ell^1$ bookkeeping and sharp product bounds along
Tao's update process, we obtain a complete conditional proof of
$\delta_m\le \mathrm{poly}(m)\,2^{-0.0504\,m}$ whose single unproven input
is a local equidistribution statement for the orbit $\{2^e \bmod 3^n\}$ in
windows of length $O(\log q)$ --- a finite-orbit quantitative form of
Furstenberg $\times 2\times 3$ rigidity that we show lies outside current
methods. For each fixed $m$ the input is decidable, and we certify the
resulting unconditional bounds at $m=64$, $128$, $256$, with the certified
rate improving toward $0.0504$. Certified lattice computations to $m=1200$
support the sharp-rate conjecture
$\delta_m=\Theta(2^{-I_0 m}m^{-c})$: the extremal frequency tracks the
cycle-resonance line $2^{k^*}\approx 3^m$ at constant deficit
$k^*=m\log_2 3-(7.0\pm 0.3)$, and the extremal amplitude is generated by a
single pinned boundary-layer bridge which already attains
$2^{-I_0m}m^{-1.6}$ with only polynomial phase loss.
\end{abstract}

\maketitle

\section{Introduction}\label{sec:intro}

The Collatz ($3x+1$) conjecture concerns the map on positive integers that
sends even $n$ to $n/2$ and odd $n$ to $3n+1$. The strongest known result,
due to Tao \cite{Tao2022}, is that almost all orbits (in logarithmic
density) attain almost bounded values. At the analytic heart of that proof
sits a single quantitative object: the $3$-adic equidistribution of the
\emph{Syracuse random variable}, the random $3$-adic integer
\[
\mathbf{Syrac}\;=\;\sum_{j\ge 1} 3^{j-1}2^{-A_j},\qquad
A_j=a_1+\dots+a_j,\quad a_i\ \text{i.i.d.},\ \Prob(a_i=a)=2^{-a},
\]
which describes the residues mod $3^m$ of the accelerated Collatz
iteration started from a random odd integer. Writing
$\e(x)=e^{2\pi i x}$, the quantity that controls everything is the maximal
nontrivial Fourier coefficient
\[
\delta_m\;=\;\max_{\xi\in\Z/3^m\Z,\ 3\nmid \xi}
\bigl|\E\,\e(\xi\,\mathbf{Syrac}/3^m)\bigr|.
\]
Tao's Proposition~1.17 proves $\delta_m\ll_A m^{-A}$ for every fixed $A$
--- superpolynomial decay --- and Section~1.4 of \cite{Tao2022} states
explicitly that an improvement to $O(\exp(-cm))$ may be possible but is
not needed there and not pursued. A survey of the citing literature
(Section~\ref{sec:cond}) shows that no later work supplies the
exponential improvement. Its value would be more than aesthetic: the rate
of $\delta_m$ is the rate-limiting constant in any \emph{effective}
version of Tao's theorem.

However, the exact decay rate of $\delta_m$ --- even its numerical value
--- appears never to have been determined. Here we determine it
numerically with certified lattice computations, identify its mechanism,
prove the structural theorems that reduce the upper-bound problem to a
single isolated arithmetic statement, and certify unconditional
exponential bounds for fixed levels $m\le 256$. Our results are of four
kinds, and we keep the four statuses --- \emph{theorem},
\emph{computational fact} (exact or certified computation),
\emph{conditional theorem}, and \emph{conjecture} --- explicitly labeled
throughout.

\subsection*{Theorems} (i) Since $2$ is a primitive root mod $3^m$, every
nontrivial frequency is a power of $2$, and the two-variable Fourier
problem collapses, with no assumptions, to a one-dimensional recursion
over the exponent lattice, with $\delta_m$ non-increasing in $m$ and the
underlying residue chain exactly stationary after $m$ steps
(Theorem~\ref{thm:onedim}). (ii) A first-passage renormalization identity
expresses $g_m$ in terms of shorter chains, with the certified kernel
bound $|F(d,r)|\le 2^{-r}\delta_{d-1}$ and an exact identification of the
phase-free region with the characteristic function of the real Syracuse
variable (Theorem~\ref{thm:renorm}). (iii) A weighted contraction lemma:
with weight $2^{\lambda u}$ in the coordinate $u=k-m\log_2 3$, the
phase-free weighted row sum equals $f(\lambda)=3^{\lambda}/(2^{1+\lambda}-1)$,
and
\[
\min_{\lambda}f(\lambda)\;=\;2^{-I_0},\qquad
I_0=\frac{1-H(\theta)}{\theta}=D\bigl(\Geom(\theta)\,\|\,\Geom(1/2)\bigr)
=0.0793186\ldots,
\]
$\theta=\log_3 2$, $H$ the binary entropy (Theorem~\ref{thm:weight}) ---
below the critical line no phase cancellation is needed to contract at
exactly the conjectured rate, so the upper-bound problem reduces to the
near-line and above-line frequencies. (iv) The real Syracuse variable is a
Kesten chain whose Kesten equation $\E[(3\cdot 2^{-a})^{\kappa}]=1$ holds
at $\kappa=1$ \emph{exactly} (the identity $3=2^2-1$), placing it at the
Cauchy-type boundary with log-periodic corrections
(Theorem~\ref{thm:kesten}). (v) An exact $2$-adic law on full periods: the
identity $\nu_2(r(K))\ge\sigma\iff |r(K-\sigma)|\le M/2^{\sigma+1}$ for
signed residues $r(K)$ of $2^K$ mod $M=3^m$ converts $2$-adic valuation
into Archimedean size, and yields that within any symmetric ball the
valuation law over a full period is geometric with explicit error
(Theorem~\ref{thm:cstat}).

\subsection*{Conditional theorem} Assembling (i)--(v) with a
capped weighted-$\ell^1$ bookkeeping (audited exactly on all $399$
computed levels, with zero violations) and sharp product bounds along
Tao's two-dimensional update process, we obtain a complete proof
architecture for
\[
\delta_m\;\le\;\mathrm{poly}(m)\cdot 2^{-0.0504\,m}
\]
--- the qualitative form of the improvement flagged in \cite{Tao2022} ---
whose \emph{single} unproven input (Hypothesis~\ref{hyp:clocal}) is a
local equidistribution statement: entries of the orbit
$\{2^e\bmod 3^n\}$ into a short symmetric interval have geometrically
decaying $2$-adic valuations, uniformly over windows of length
$O(\log q)$, $q=3^n$. We show (Section~\ref{sec:cond}) that this
statement is a finite-orbit quantitative form of Furstenberg's
$\times 2\times 3$ phenomenon \cite{Furstenberg1967}, that the effective
results of Bourgain--Lindenstrauss--Michel--Venkatesh \cite{BLMV2009} do
not reach it, and that the window length $O(\log q)$ is exponentially
below the reach of the known exponential-sum technology; it is, however,
decidable in time $O(m^4)$ for each fixed $m$.

\subsection*{Computational facts} Certified exact-arithmetic computations
give: full-frequency scans for $m\le 15$; rescaled window computations to
$m=1200$ showing that the extremal frequency exponent satisfies
$k^*=m\log_2 3-(7.0\pm0.3)$ with the deficit \emph{constant} over
$m=200$--$1200$, and a decay rate in $[0.0790, 0.0800]$ bits per level
approaching $I_0$ one-sidedly; a boundary-layer structure theory
(Section~\ref{sec:mech}); and unconditional certified instances
\[
\delta_{64}\le C\,2^{-0.0350\cdot 64},\qquad
\delta_{128}\le C\,2^{-0.0406\cdot 128},\qquad
\delta_{256}\le C\,2^{-0.0457\cdot 256},
\]
with all band constraints closing at positive slack and the certified
rate increasing toward $0.0504$ (Section~\ref{sec:instances}).

\subsection*{Conjecture} The sharp rate
(Conjecture~\ref{conj:c1}): $\delta_m=\Theta(2^{-I_0m}\,m^{-c})$.
Equivalently, the $3$-adic mixing bottleneck of Syracuse is the
cycle-resonance line $2^{A}\approx 3^{k}$, and the mixing rate coincides,
after clock conversion, with the classical stopping-time density exponent
$1-H(\theta)$: one constant read in three clocks.

\subsection*{Method} This paper practices experimental mathematics with
pre-registered decision criteria: for each major numerical discrimination
we declared acceptance/rejection thresholds before computing, and we
report the outcomes --- including a rejected candidate rate --- in
Appendix~\ref{app:data}, which doubles as a calibration log. All certified
quantities use exact integer or interval-safe arithmetic; scripts and
frozen data accompany the paper. A companion paper \cite{P1} studies the
$2$-adic (certificate-complexity) side of the same constant.

\section{Setup and the one-dimensional reduction}\label{sec:onedim}

Fix $m\ge 1$ and let $w_0$ be arbitrary in $\Z/3^m\Z$. The Syracuse
residue chain is
\begin{equation}\label{eq:chain}
w_t\;=\;(3w_{t-1}+1)\cdot 2^{-a_t} \bmod 3^m,\qquad
a_t\ \text{i.i.d.},\ \Prob(a_t=a)=2^{-a}\ (a\ge 1),
\end{equation}
where $2^{-1}$ denotes the inverse of $2$ mod $3^m$. This is the law of
the residues mod $3^m$ of the accelerated Collatz map
$n\mapsto (3n+1)/2^{\nu_2(3n+1)}$ along the odd iterates, in the sense
made precise in \cite[\S1]{Tao2022}; the stationary law of
\eqref{eq:chain} is the law of $\mathbf{Syrac}$ mod $3^m$.

\begin{theorem}[one-dimensional reduction]\label{thm:onedim}
\begin{enumerate}
\item[(a)] $2$ is a primitive root mod $3^m$ for every $m\ge 1$:
  $\ord(2 \bmod 3^m)=2\cdot 3^{m-1}=|(\Z/3^m\Z)^{\times}|$.
\item[(b)] Consequently every $\xi$ with $3\nmid\xi$ is $\pm$ a power of
  $2$ mod $3^m$, and with
  $g_m(k):=\E\,\e\!\bigl(2^{k}\,\mathbf{Syrac}/3^m\bigr)$ (exponents read
  mod $\ord(2)$) one has $\delta_m=\max_k |g_m(k)|$ and the exact lattice
  recursion
  \begin{equation}\label{eq:rec}
  g_m(k)\;=\;\sum_{a\ge 1}2^{-a}\,
  \e\!\Bigl(\tfrac{2^{\,k-a} \bmod 3^m}{3^m}\Bigr)\,g_{m-1}(k-a).
  \end{equation}
\item[(c)] The chain \eqref{eq:chain} is \emph{exactly} stationary after
  $m$ steps: for $t\ge m$ the law of $w_t$ equals the law of
  $\mathbf{Syrac}$ mod $3^m$, for every initial $w_0$.
\item[(d)] $\delta_m\le\delta_{m-1}$ for all $m\ge 2$.
\end{enumerate}
\end{theorem}

\begin{proof}
(a) From $2^2=1+3$ and the lifting-the-exponent identity
$(1+3^s u)^3\equiv 1+3^{s+1}u \pmod{3^{s+2}}$ one gets
$\nu_3(2^{2\cdot 3^{t}}-1)=t+1$ exactly, so the order of $2$ mod $3^m$
divides $2\cdot 3^{m-1}$ and is divisible by $3^{m-1}$; and
$2^{3^j}\equiv (-1)^{3^j}=-1\not\equiv 1 \pmod 3$ shows the order is
even. Hence $\ord(2)=2\cdot 3^{m-1}$.

(b) By (a) the cyclic group generated by $2$ is the whole unit group, so
the powers of $2$ exhaust all $\xi$ with $3\nmid\xi$, and the maximum
defining $\delta_m$ is attained on the exponent lattice. The recursion follows by
conditioning \eqref{eq:chain} on the last step: with $\xi=2^k$,
\[
\e\!\bigl(\xi(3w+1)2^{-a}/3^m\bigr)
=\e\!\bigl(2^{k-a}\cdot 3w/3^m\bigr)\,\e\!\bigl(2^{k-a}/3^m\bigr)
=\e\!\bigl(2^{k-a}w/3^{m-1}\bigr)\,\e\!\bigl(2^{k-a}/3^m\bigr),
\]
and taking expectations over $w\sim\mathbf{Syrac}\bmod 3^{m-1}$ and $a$
gives \eqref{eq:rec}.

(c) By induction on $j$: $w_t\bmod 3^j$ depends only on
$(a_{t-j+1},\dots,a_t)$, because $(3w+1)\bmod 3^{j}$ depends only on
$w\bmod 3^{j-1}$ and multiplication by the unit $2^{-a}$ preserves this.
For $j=m$, $t\ge m$, the law of $w_t$ is a fixed function of the last $m$
steps, independent of $w_0$; it is therefore the unique stationary law.

(d) Bound \eqref{eq:rec} by the triangle inequality:
$|g_m(k)|\le\sum_a 2^{-a}|g_{m-1}(k-a)|\le\delta_{m-1}$, using that every
frequency $2^{k-a}$ is again a unit mod $3^{m-1}$.
\end{proof}

\begin{remark}
Theorem~\ref{thm:onedim}(b) upgrades to a theorem what would otherwise be
a heuristic restriction to special frequencies: \emph{every} nontrivial
frequency lives on the exponent lattice, so the recursion
\eqref{eq:rec} loses nothing. Full-frequency exhaustive scans for
$m\le 15$ confirm that the maximizer is always a power of $2$ in a
narrow window of exponents (Section~\ref{sec:numerics}).
\end{remark}

\section{First-passage renormalization}\label{sec:renorm}

Let $Y_m=\sum_{j=1}^m 3^{j-1}2^{-A_j}$ denote the \emph{real} Syracuse
variable (the same series read in $\R$), and write $x=2^k/3^m$.

\begin{theorem}[renormalization identity]\label{thm:renorm}
For all $m,k$:
\begin{enumerate}
\item[(a)] (phase-free region) On the event $\{A_m\le k\}$ the $3$-adic
  phase coincides exactly with the real one:
  $\e(2^k\,\mathbf{Syrac}_m/3^m)=\e(2^k Y_m/3^m)$ pathwise, where
  $\mathbf{Syrac}_m$ is the level-$m$ truncation. Hence
  $\E[\e(2^k\mathbf{Syrac}_m/3^m);A_m\le k]=\E[\e(xY_m);A_m\le k]$.
\item[(b)] (first-passage decomposition) With
  $\mathrm{amp}_j(A)=\E[\e(xY_j); A_j=A]$ and
  \[
  F(d,r)\;=\;\sum_{a>r}2^{-a}\,
  \e\!\Bigl(\tfrac{2^{\,r-a}\bmod 3^d}{3^d}\Bigr)\,V_{d-1}(2^{\,r-a}),
  \qquad V_{d-1}(\xi)=\E\,\e(\xi\,\mathbf{Syrac}_{d-1}/3^{d-1}),
  \]
  one has the exact identity
  \[
  g_m(k)\;=\;\E[\e(xY_m);A_m\le k]\;+\;
  \sum_{j<m}\ \sum_{A\le k}\mathrm{amp}_j(A)\,F(m-j,\,k-A),
  \]
  obtained by splitting at the first step where the valuation sum
  exceeds $k$.
\item[(c)] (kernel bound) $|F(d,r)|\le 2^{-r}\,\delta_{d-1}$.
\end{enumerate}
\end{theorem}

\begin{proof}
(a) On $\{A_m\le k\}$ every exponent $k-A_j$ ($j\le m$) is non-negative,
so $2^{k}\cdot(2^{-A_j}\bmod 3^m)\equiv 2^{k-A_j}$ mod $3^m$ with
$2^{k-A_j}$ a genuine integer, and
$\e\bigl(\sum_j 3^{j-1}2^{k-A_j}/3^m\bigr)$ is simultaneously the
$3$-adic and the real phase.

(b) Decompose according to the first index $j+1$ with $A_{j+1}>k$
(if none, we are in the event of (a)). Conditioning on $A_j=A\le k$ and
on the overshooting step $a_{j+1}=a>k-A$, the remaining $m-j-1$ steps
form a fresh chain at level $m-j-1$, and collecting phases gives the
stated kernel with $r=k-A$; the identity is exact. (It was additionally
verified numerically to relative error $10^{-14}$ across levels.)

(c) $|F(d,r)|\le\sum_{a>r}2^{-a}\,|V_{d-1}(2^{r-a})|
\le \delta_{d-1}\sum_{a>r}2^{-a}=2^{-r}\delta_{d-1}$, again because the
frequencies are units.
\end{proof}

\begin{remark}[why absolute values do not suffice]\label{rem:noabs}
It is tempting to iterate Theorem~\ref{thm:renorm}(b)--(c) with absolute
values throughout. This provably fails to close: the phase-free main term
must then be bounded by the large-deviation probability
$\Prob(A_m\le k)$, and near the critical line the inductive inequality
loses a constant factor per level that destroys any exponent. The failure
is structural --- on the good event the final phase still rotates by at
least $1/3$ of a turn, so positivity-based lower bounds are equally
impossible, and the modulus of $g_m$ is carried by cancellation among
crossing terms (Section~\ref{sec:mech}). This motivates the weighted
approach of the next section.
\end{remark}

\section{Weighted contraction and the reduction of the upper bound}
\label{sec:weight}

Write $\alpha=\log_2 3$ and use the line coordinate $u=k-m\alpha$.

\begin{theorem}[weighted contraction]\label{thm:weight}
For $\lambda>-1$ define the phase-free weighted row sum of the recursion
\eqref{eq:rec} with weight $2^{\lambda u}$:
\[
f(\lambda)\;=\;\sum_{a\ge 1}2^{-a}\,2^{\lambda(\alpha-a)}
\;=\;\frac{3^{\lambda}}{2^{1+\lambda}-1}.
\]
Then $f$ is strictly log-convex on $(-1,\infty)$ and
\[
\min_{\lambda}f(\lambda)\;=\;2^{-I_0},\qquad
I_0=\frac{1-H(\theta)}{\theta},\qquad
\text{attained at } 2^{1+\lambda^*}=\frac{\alpha}{\alpha-1},\quad
\lambda^*=0.4380326\ldots
\]
\end{theorem}

Applying the weight to \eqref{eq:rec} and bounding phases by $1$
wherever the weight decays, one obtains for the weighted norm
$\|g_m\|_{\lambda}=\sup_k 2^{\lambda u}|g_m(k)|$ the one-step inequality
$\|g_m\|_{\lambda}\le f(\lambda)\|g_{m-1}\|_{\lambda}+(\text{boundary
terms supported near and above } u=0)$: with no phase input at all, the
below-line region contracts at exactly the conjectured rate, and the
upper-bound problem for $\delta_m$ reduces to the control of frequencies
with $u\gtrsim 0$. The precise audited form of this reduction is the
capped functional described after the proof.

\begin{proof}
The row sum is a geometric series: $\sum_a 2^{-a(1+\lambda)}=1/(2^{1+\lambda}-1)$,
times $2^{\lambda\alpha}=3^{\lambda}$. For the minimum, substitute
$t=2^{1+\lambda}>1/2$: then $3^{\lambda}=t^{\alpha}/3$ and
$f=t^{\alpha}/(3(t-1))$, with critical point $\alpha(t-1)=t$, i.e.\
$t^*=\alpha/(\alpha-1)$, and
\[
\min f=\frac{(t^*)^{\alpha}}{3(t^*-1)}
=\frac{\alpha^{\alpha}}{3\,(\alpha-1)^{\alpha-1}} .
\]
Taking $-\log_2$ and using $\log_2 3=\alpha$:
$-\log_2\min f=\alpha-\alpha\log_2\alpha+(\alpha-1)\log_2(\alpha-1)$.
On the other hand, with $\theta=1/\alpha$,
\[
\frac{1-H(\theta)}{\theta}
=\alpha\Bigl(1-\tfrac1\alpha\log_2\alpha
-\bigl(1-\tfrac1\alpha\bigr)\log_2\tfrac{\alpha}{\alpha-1}\Bigr)
=\alpha-\alpha\log_2\alpha+(\alpha-1)\log_2(\alpha-1),
\]
which is the same expression. (The identity was additionally checked to
$50$ digits.) Log-convexity holds because $f(\lambda)$ is the moment
generating function $\E\,2^{\lambda(\alpha-a)}$ of the step variable.
\end{proof}

Two consequences organize everything that follows. First, $I_0$ is
simultaneously a Kullback--Leibler divergence:
\begin{equation}\label{eq:kl}
I_0\;=\;D\bigl(\Geom(\theta)\,\|\,\Geom(1/2)\bigr),\qquad
\theta\, I_0=1-H(\theta),
\end{equation}
the second equality being a one-line identity; the tilted step law at the
optimum is exactly $\Geom(\theta)$, whose mean step is $1/\theta=\alpha$
--- the drift that tracks the critical line. Second, the reduction can be
run with a \emph{free} weight exponent. For $\mu<0$ define
\begin{equation}\label{eq:rmu}
r(\mu)\;=\;\mu\,\alpha+\log_2\!\bigl(2^{\,1-\mu}-1\bigr),
\end{equation}
the contraction rate of the capped weighted-$\ell^1$ functional
$N_{\mu}(m)=\sum_k \min\bigl(2^{\mu u},\,1\bigr)\,|g_m(k)|$; one checks
$\max_{\mu}r(\mu)=I_0$ at $\mu=-\lambda^*$, and smaller $|\mu|$ trades
rate for weaker demands on the above-line frequencies. The functional
$N_\mu$ has no above-line escape route: the weight decays upward, and,
unlike the sup norm, the $\ell^1$ sum sees every frequency.

\begin{compfact}[audited bookkeeping]\label{cf:audit}
The one-step inequality $N_{\mu}(m)\le 2^{-r(\mu)}N_{\mu}(m-1)$ with
$\mu=-0.16$, $r(\mu)=0.0504$, was audited against the exact lattice data
on all $399$ computed levels: zero violations, with no injection terms
required within the window. The remaining bookkeeping (seam and
wrap-around of the exponent circle) is exactly the constraint system
used in the certified instances of Section~\ref{sec:instances}.
\end{compfact}

\section{Kesten structure of the real variable}\label{sec:kesten}

\begin{theorem}\label{thm:kesten}
The real Syracuse variable $Y=\sum_{j\ge1}3^{j-1}2^{-A_j}$ satisfies the
stochastic fixed-point equation
$Y\overset{d}{=}3\cdot 2^{-a}\,Y+2^{-a}$ ($a\sim\Geom(1/2)$ independent
of $Y$ on the right), and its Kesten equation
$\E[(3\cdot 2^{-a})^{\kappa}]=1$ holds \emph{exactly} at $\kappa=1$:
\[
\E\bigl[3\cdot 2^{-a}\bigr]
=\sum_{a\ge 1}2^{-a}\cdot 3\cdot 2^{-a}
=3\sum_{a\ge 1}4^{-a}=3\cdot\tfrac13=1 .
\]
The multiplier $\log(3\cdot 2^{-a})=\log 3-a\log 2$ is
arithmetically lattice, and the Goldie conditions hold; hence
$\Prob(Y>t)=h(\log t)/t\,(1+o(1))$ with $h$ log-periodic, $\E Y=\infty$,
and the tail exponent is exactly $1$.
\end{theorem}

\begin{proof}
The fixed-point equation restates the definition after one step. The
Kesten computation at $\kappa=1$ is the displayed line ---
$3^{\kappa}=2^{\kappa+1}-1$ has the integer solution $\kappa=1$ because
$3=2^2-1$. The moment and lattice conditions of implicit renewal theory
\cite{Kesten1973,Goldie1991} are immediate ($\E[(3\cdot2^{-a})\log^+(3\cdot2^{-a})]<\infty$),
giving the tail with log-periodic prefactor in the lattice case.
\end{proof}

This places Syracuse at the \emph{boundary} Cauchy-type case: power
corrections in $1/m$ and log-periodic constants are forced, which is
exactly what the numerics of Section~\ref{sec:numerics} display, and the
polynomial factor in Conjecture~\ref{conj:c1} should be regarded as this
mechanism's signature rather than as numerical noise.

\section{An exact 2-adic law on full periods}\label{sec:cstat}

The conditional route of Section~\ref{sec:cond} needs one probabilistic
input about the deterministic orbit $\{2^K \bmod 3^m\}$: entries into a
short symmetric interval should have geometrically distributed $2$-adic
valuations. On \emph{full periods} this is a theorem, by an elementary
identity that converts the $2$-adic condition into an Archimedean one.

For odd $M\ge 3$ let $r(K)\in(-M/2,M/2]$ denote the signed residue of
$2^K$ mod $M$.

\begin{lemma}[valuation--size identity]\label{lem:val}
For all $K$ and $\sigma\ge 0$:
\[
\nu_2\bigl(r(K)\bigr)\ \ge\ \sigma
\qquad\Longleftrightarrow\qquad
\bigl|r(K-\sigma)\bigr|\ \le\ \frac{M}{2^{\sigma+1}} .
\]
\end{lemma}

\begin{proof}
($\Rightarrow$) If $2^{\sigma}\mid r(K)$ then $r(K)/2^{\sigma}$ is an
integer congruent to $2^{K-\sigma}$ mod $M$ with absolute value at most
$M/2^{\sigma+1}<M/2$, hence it \emph{is} the signed residue $r(K-\sigma)$.
($\Leftarrow$) If $|r(K-\sigma)|\le M/2^{\sigma+1}$ then
$2^{\sigma}r(K-\sigma)$ is congruent to $2^{K}$ mod $M$ with absolute
value at most $M/2$, hence equals $r(K)$, which is thus divisible by
$2^{\sigma}$.
\end{proof}

\begin{theorem}[geometric law on cycles]\label{thm:cstat}
Let $M=3^m$ and let $K$ range over a full period of the orbit
$\{2^K \bmod 3^m\}$ (length $2\cdot 3^{m-1}$, all units, by
Theorem~\ref{thm:onedim}(a)). Then for every $\varepsilon\in(0,1/2]$ and
every $\sigma\ge 0$,
\[
\Prob_K\Bigl(\nu_2(r(K))\ge\sigma\ \Bigm|\ |r(K)|\le\varepsilon 3^m\Bigr)
\;=\;2^{-\sigma}\;+\;O\!\Bigl(\frac{2^{\sigma}}{\varepsilon\,3^{m}}\Bigr),
\]
i.e.\ inside any symmetric ball the valuation law over a full period is
geometric of ratio $1/2$, with constant $C=1$, up to the stated counting
error.
\end{theorem}

\begin{proof}
By Lemma~\ref{lem:val} (applied once more to intersect with the ball),
$\{\,\nu_2(r(K))\ge\sigma,\ |r(K)|\le\varepsilon 3^m\}$ is, under
$K\mapsto K-\sigma$, in bijection with
$\{\,|r(K')|\le\varepsilon 3^m/2^{\sigma}\}$. Since $2$ is a primitive
root, over a full period the map $K\mapsto 2^K$ visits every unit
exactly once, so the count of $K'$ with $|r(K')|\le x$ equals the number
of nonzero integers prime to $3$ in $[-x,x]$, which is
$\tfrac{4}{3}x+O(1)$. Taking the ratio of the counts
at $x=\varepsilon 3^m/2^{\sigma}$ and $x=\varepsilon 3^m$ gives
$2^{-\sigma}(1+O(2^{\sigma}/\varepsilon 3^m))$.
\end{proof}

\begin{remark}
An independent derivation via Palm sampling of ball entries and a
conductance comparison gives the same conclusion without the bijection.
Exhaustive verification at $m=12$ over the full period and
$\sigma\le 5$ found zero mismatches of Lemma~\ref{lem:val}, and the
observed frequencies $0.4999,\,0.2501,\,0.1250,\,0.0624$ match the
geometric law. We also note a caveat used below: value chains generated
by non-units (multiples of $9$) live on leaves of effective modulus
$3^{m-2j+2}$, and Lemma~\ref{lem:val} is applied at the effective scale.
\end{remark}

\section{Certified numerics to $m=1200$}\label{sec:numerics}

All computations use the recursion \eqref{eq:rec} on the exponent
lattice, in exact or error-tracked arithmetic; scripts and data
accompany the paper.

\begin{compfact}[extremal frequency and rate]\label{cf:numerics}
\begin{enumerate}
\item[(a)] Full-frequency scans for $m\le 15$ (all $2\cdot 3^{m-1}$
  frequencies): the maximizer of $|g_m|$ is always a power of $2$ with
  exponent in a narrow window at the critical line; the window
  computation reproduces the true $\delta_m$ to ratio $1.0000$.
\item[(b)] Rescaled lattice computations for $m\le 1200$: the extremal
  exponent satisfies $k^*=m\log_2 3-(7.0\pm 0.3)$, with the deficit
  constant over $m=200$--$1200$; the local decay rate of the window
  maximum lies in $[0.0790,\,0.0800]$ bits/level and approaches
  $I_0=0.07932$ from above, consistent with a polynomial correction.
\item[(c)] The one-sided nature of (b) matches the theory: weighted
  contraction (Theorem~\ref{thm:weight}) can only prove rate $\ge I_0-\epsilon$
  statements for the certified functionals, and the tilted optimum
  behaves as $\Geom(\theta)$.
\item[(d)] The above-line profile is monotone and steep: the sup over
  $u\in(0,2]$ is $4.6\%$ of the ridge value, $9\times 10^{-5}$ over
  $u\in(10,20]$, $4\times 10^{-13}$ over $u\in(40,80]$ (computed in a
  truncation-free extended window at $m=100,200$): there are no deep
  resonances above the line.
\end{enumerate}
\end{compfact}

The caveat that a window computation might miss a distant frequency
among the $\sim 10^{286}$ at $m=600$ is answered structurally: the
$\ell^4$ functional $S_4'(m)=\sum_{3\nmid\xi}|\E\e(\xi X/3^m)|^4$ is a
\emph{computable global scalar} --- an FFT sees every frequency --- with
the collision representation
$S_4'(m)=3^m\,\Prob(X_1+X_2\equiv X_3+X_4 \bmod 3^m)$; exact computation
for $m\le 13$ shows $Q(m)=S_4'/\delta_m^4$ saturating at $\approx 18$, so
the fourth moment is concentrated on the extremal ridge and
$\delta_m=(S_4'/Q)^{1/4}$. The four-fold weighted row sum factorizes
exactly, $\min_{\lambda}[3^{-\lambda}/(2^{1-\lambda}-1)]^4=2^{-4I_0}$
(the joint four-walk representation is essential: Jensen on the marginal
yields only $2^{-I_0}$), and the windowed $\ell^4$ per-level rate
converges one-sidedly to $4I_0=0.31727$ ($0.403\to0.332$ at $m=400$).

\section{The boundary-layer mechanism}\label{sec:mech}

Where does the extremal amplitude live? The following structure
statements are certified computations (with pre-registered criteria;
outcomes, including one rejected hypothesis, in
Appendix~\ref{app:data}).

\begin{compfact}[boundary layer]\label{cf:mech}
\begin{enumerate}
\item[(a)] (main-term failure) The phase-free main term
  $\E[\e(xY_m);A_m\le k^*]$ decays at rate $I_0+0.0093$ --- strictly
  faster than $\delta_m$; positivity of the good event is destroyed
  exponentially, and crossing terms carry $80$--$257$ times the main
  term. Any proof of the lower bound must be an oscillatory one.
\item[(b)] (corner localization) In the first-passage decomposition of
  Theorem~\ref{thm:renorm}, $70\%$ of the contribution mass to
  $g_m(k^*)$ sits in $9$--$18$ cells around the corner
  $r=k^*-A=0$, $d=m-j=5$, with $d$ \emph{exactly constant} over
  $m=24$--$100$, and geometric decay $2^{-1.5r}$ in $r$: the extremal
  amplitude is generated by paths whose valuation sum first reaches the
  critical value $O(1)$ levels before the end --- a single pinned
  boundary-layer bridge.
\item[(c)] (pinned bridge attains the rate) The single pinned amplitude
  obeys $|\mathrm{amp}_{m-5}(k^*)|=C\,2^{-a m}m^{-1.59}$ with
  $a=0.07914\in[I_0-0.001,\,I_0+0.002]$ for $m\le 400$: no exponential
  phase cancellation occurs in the bridge. Phase retention decays
  polynomially, $\sim m^{-1.06}$, distributed across scales: the octave
  factors of the retention profile $K(s,m)$ (which is $m$-independent
  for $s\ll m$) lie in $[0.62,0.95]$ for $s\in[4,256]$ at $m=640$ ---
  uniformly bounded away from $0$, with no catastrophic octave.
\item[(d)] (decorrelation length) The autocorrelation of $g_m$ along the
  exponent lattice has universal profile with correlation length
  $\approx 8$ exponent units (matching the ridge width and the deficit
  $\approx 7$), measured identically at $m=200$ and $400$; this is the
  mechanism separating the joint rate $4I_0$ from the marginal $I_0$ in
  the $\ell^4$ functional.
\end{enumerate}
\end{compfact}

Summarizing the mechanism: $\delta_m$ is a boundary-layer phenomenon on
the cycle-resonance line $2^{A}\approx 3^{k}$ --- the same Diophantine
line that governs the nontrivial-cycle problem --- with amplitude carried
by a single pinned crossing and only polynomial phase loss. The two
remaining asymptotic lemmas are stated in Section~\ref{sec:conj}.

\section{Toward unconditional exponential decay: the conditional theorem}
\label{sec:cond}

\subsection{Literature status}
Proposition~1.17 of \cite{Tao2022} gives $\delta_m\ll_A m^{-A}$,
uniformly over $3\nmid\xi$; the end of \cite[\S1.4]{Tao2022} states the
possible improvement to $O(\exp(-cm))$ and explicitly does not pursue
it. A scan of the citing literature ($144$ citing works at title level,
plus arXiv metadata searches for the Syracuse characteristic function)
found no paper supplying the exponential improvement. The mechanism of
\cite[\S7]{Tao2022} is a two-dimensional renewal process whose estimate
is driven by the number of ``white points'' along the walk; the
$m^{-A}$ loss comes from the largest black triangle, and an exponential
bound needs a \emph{linear} lower bound on white points --- which is
what the structure below supplies, conditionally.

\subsection{The proven components}
Along a wrapped segment of the walk, the relevant object is the value
chain $v(e)=\mathrm{signed}(3^{2j-2}\,2^{e} \bmod 3^d)$ at effective
scale $d$, with the black set $B_{\varepsilon}=\{|v|\le\varepsilon 3^d\}$
(Tao's triangles) and its dyadic subsets
$D_s=\{v\in B_{\varepsilon}:\nu_2(v)\ge s\}$.

\begin{lemma}[doubling without reduction]\label{lem:r1a}
If $|v(e)|\le 3^d/4$ then $v(e+1)=2v(e)$ exactly. Consequently each
maximal black run is the doubling chain of a unique odd seed $w$
(the odd part of $v$), $\{v=2^tw:\ 0\le t\le\log_2(\varepsilon 3^d/|w|)\}$,
and chains of distinct seeds are disjoint; the seed $w=3^{2j-2}$ (when
below $\varepsilon 3^d$) is the literal small-integer main triangle.
\end{lemma}

\begin{lemma}[pair spacing]\label{lem:pair}
Two seeds of depth $h$ (i.e.\ chains reaching $\nu_2\ge h$) in the same
row are separated by at least $2\alpha h-\log_2((1+\varepsilon)/\varepsilon)$
positions. (Cross-ratio argument on $t\neq 0$ intersections; this is the
explicit quantitative form of the separation of large triangles.)
\end{lemma}

\begin{lemma}[$2$-adic exhaustion, probabilistic form]\label{lem:exhaust}
Within a black chain, continuation requires $b_j\le\nu_2(v)$ and the
budget $\nu_2$ decreases by $b$ at each step; deterministic exit is
forced for $b\le\floor{\log_2(1/2\varepsilon)}$, while larger $b$ can
re-seed only by direct hits of measure $O(\varepsilon 2^{b}/M)$ per
step (measured: $0.2\%$ of transitions). Hence the expected black run
per entry is $O(1)$ with geometric tail, \emph{given} the entry-valuation
hypothesis below.
\end{lemma}

Lemma~\ref{lem:r1a} is proved by the displayed inequality
($|2v|\le 3^d/2$ precludes reduction); Lemma~\ref{lem:pair} is a
four-line cross-ratio argument; Lemma~\ref{lem:exhaust} combines the
budget bound with the measured re-seeding rate (the deterministic-exit
claim for all $b$ is \emph{false} --- an explicit counterexample at
$m=7$, $v=-13$, $b=7$ exists --- which is why the run-length control is
routed through the entry-valuation hypothesis).

\subsection{The single remaining input}

\begin{hypothesis}[local entry-valuation bound; ``Lemma C-local'']
\label{hyp:clocal}
There is an absolute constant $C$ such that, along wrapped segments of
length $O(m)$ of the orbit $\{2^e \bmod 3^d\}$ at every effective scale
$d\le m$,
\[
\sum_t \Prob\bigl(v_{t-1}\notin B,\ v_t\in D_s\bigr)
\;\le\; C\,2^{-s}\,\sum_t \Prob\bigl(v_{t-1}\notin B,\ v_t\in B\bigr)
\qquad (s\ge 0):
\]
fresh entries into the ball land on the depth-$s$ subset with at most
geometric proportion.
\end{hypothesis}

Theorem~\ref{thm:cstat} proves exactly this with $C=1$ when the sums
run over \emph{full periods}; the hypothesis is its restriction to
windows of length $O(m)=O(\log q)$, $q=3^d$. Measured behaviour: in the
deep bands the entry valuations match $\Geom(1/2)$ to three digits
($0.482/0.272/0.129/0.067$); in the shallow shoulder band the tail is
enhanced by a bounded factor ($\le 5$ at the measured depths), which is
harmless since Hypothesis~\ref{hyp:clocal} permits any absolute $C$.

\subsection{Classification of the hypothesis}
Hypothesis~\ref{hyp:clocal} is a \emph{finite-orbit, quantitative} form
of the Furstenberg $\times 2\times 3$ phenomenon
\cite{Furstenberg1967}: it asserts joint $2$-adic/$3$-adic
non-concentration for the single multiplicative orbit of $2$ modulo
powers of $3$, in windows of length logarithmic in the modulus. The
effective $\times a\times b$ results of \cite{BLMV2009} concern
semigroup orbits and do not reach a single $\times 2$ orbit; the
strongest exponential-sum technology for power residues (Postnikov's
character formula with Vinogradov/Korobov estimates, and power-modulus
refinements) becomes nontrivial only for windows of length
$\ge \exp(c(\log q)^{2/3+\epsilon})$, exponentially longer than the
$O(\log q)$ required; additive-combinatorial methods require set sizes
$q^{\delta}$. The hypothesis is also adjacent to (but strictly weaker
in its averaged form than) the Erd\H{o}s ternary-digit problem for
$2^n$ \cite{Erdos1979,Lagarias2009}. We therefore regard it as
precisely isolated but genuinely open; it is, however, decidable for
each fixed $m$, which is what Section~\ref{sec:instances} exploits.

\subsection{The conditional theorem}

\begin{theorem}[conditional exponential decay]\label{thm:conditional}
Assume Hypothesis~\ref{hyp:clocal}. Then the sharpened product bounds
along Tao's update process (all-$b$ dynamic program) admit uniform band
rates $a(\delta)$ --- $\delta=|k-m\alpha|/m$ the relative distance to
the line --- that close the capped weighted-$\ell^1$ constraint system
of Section~\ref{sec:weight} at $\mu=-0.16$ with positive slack in every
band. Consequently
\[
\delta_m\;\le\;C\,\mathrm{poly}(m)\;2^{-r m},\qquad r=0.0504\ \text{bits/level}.
\]
\end{theorem}

\begin{proof}[Proof architecture]
Hypothesis~\ref{hyp:clocal}, with Lemmas~\ref{lem:r1a}--\ref{lem:exhaust},
gives expected black-run length $O(1)$ per ball entry with geometric
tails along every wrapped segment; hence the number of white points is
linear in the segment length, at every effective scale, and the octave
sum converges (removing the logarithmic divergence responsible for
Tao's $m^{-A}$). This yields the uniform lower bounds on the band rates
$a(\delta)$ computed by the all-$b$ product DP; the measured rates
(stable in $m$ at $240$ and $320$) are: ridge $0.060$--$0.078\le I_0$
(consistency check), wrap side $0.094$--$0.17$, below-line
$0.0633/0.0810/0.108$ at $\delta=0.05/0.15/0.30$, seam $\approx 0.33$.
Optimizing the free weight over $\mu$ against the constraint
$r(\mu)+|\mu|\delta\le a(\delta)$ for all bands outside the critical
window yields $\mu=-0.16$ with slack $+0.005$--$0.030$ in every band and
rate $r(\mu)=0.0504$ by \eqref{eq:rmu}. The capped functional's one-step
contraction at this rate is the audited inequality of
Computational Fact~\ref{cf:audit}; the seam and circle-wraparound
bookkeeping close with the seam rate $0.33\gg r$. Summing the one-step
inequalities gives the claim.
\end{proof}

\begin{remark}
The division of labor is exact: every step of the architecture is either
proved (Theorems~\ref{thm:onedim}--\ref{thm:cstat},
Lemmas~\ref{lem:r1a}--\ref{lem:pair}), audited exactly
(Computational Fact~\ref{cf:audit}), or measured with stated margins
(the band rates); Hypothesis~\ref{hyp:clocal} is the sole unproven
input, and Theorem~\ref{thm:cstat} proves its full-period form. We do
not currently see how to prove the windowed form, and we flag it as the
precise obstruction --- in our view a more tractable target than the
Collatz conjecture itself, but of the $\times 2\times 3$ family.
\end{remark}

\section{Unconditional certified instances}\label{sec:instances}

For each fixed $m$, Hypothesis~\ref{hyp:clocal} is not needed: the band
inequalities are finitely many verifiable statements.

\begin{theorem}[finite certification]\label{thm:instances}
For each fixed $m$, the inequality
$\delta_m\le C\,2^{-r m}$, with $r=r(\mu)$ as in \eqref{eq:rmu} and the
constraint system of Theorem~\ref{thm:conditional} evaluated on all
window frequencies, is decidable in time $O(m^4)$. The certification
succeeds at:
\[
\begin{array}{llll}
m=64: & 207\ \text{frequencies}, & \mu=-0.10,\ r=0.0350, &
\text{min slack } +0.0078;\\
m=128: & 313\ \text{frequencies}, & \mu=-0.12,\ r=0.0406, &
\text{min slack } +0.0212;\\
m=256: & 525\ \text{frequencies}, & \mu=-0.14,\ r=0.0457, &
\text{min slack } +0.0102.
\end{array}
\]
In particular $\delta_{256}\le C\,2^{-11.7}$ unconditionally, and the
certified rate increases toward the conditional rate $0.0504$
(band-extraction value at $m=240$), with the binding constraints
stabilizing in the shoulder band $u\in[-13,-9]$ --- the same band the
asymptotic analysis identifies as the precision zone.
\end{theorem}

The computations were performed in floating point with well-conditioned
quantities and stated slacks at least $+0.008$; replacing them by outward
interval arithmetic is mechanical and changes no conclusion at the
stated precision. We regard the growth of the certified rate
($0.0350\to 0.0406\to 0.0457$, one doubling at a time) as strong
structural corroboration of Theorem~\ref{thm:conditional}.

\section{The sharp-rate conjecture and the three clocks}\label{sec:conj}

\begin{conjecture}[sharp rate]\label{conj:c1}
$\delta_m\;=\;\Theta\bigl(2^{-I_0 m}\,m^{-c}\bigr)$ with
$I_0=(1-H(\theta))/\theta=0.0793186\ldots$ and some $c>0$
(numerically $c\approx 1.6$ for the pinned bridge, $1.1$ for weighted
norms).
\end{conjecture}

Two asymptotic lemmas would upgrade the numerical support to a proof of
the corresponding bounds, and we state them precisely:
\[
(\mathrm{L}\alpha)\ \ \sup_{k\ge m\alpha+O(1)}|g_m(k)|\le 2^{-I_0m+o(m)}
\qquad\quad
(\mathrm{L}\beta)\ \ \bigl|\E_{\mathrm{tilt}}\,\e(\Phi_m)\bigr|\ \ge\ 1/\mathrm{poly}(m)
\]
--- above-line control for the upper bound, and pinned-bridge phase
non-vanishing (for the bridge pinned at $A_{m-5}=k^*$) for the lower
bound; Section~\ref{sec:mech} locates both. A structural warning is
proved along the way: at the extremal frequency the ridge phases are
$1+O(\rho^m)$, so no finite-depth expansion of the recursion can prove
either lemma --- the decay of $\delta_m$ is an unbounded-memory effect,
the Fourier-side analogue of the finite-window obstructions of
\cite{P1}.

\subsection*{Three clocks}
The constant $I_0$ is one number read in three clocks: (a) the
stopping-time/certificate density exponent $1-H(\theta)$ per binary
digit \cite{Terras1976,P1}; (b) the KL anomaly rate
$D(\Geom(\theta)\|\Geom(1/2))$ per odd step, forced on every
non-descending window by the deterministic rigidity identity; (c) the
$3$-adic mixing rate per ternary level (Conjecture~\ref{conj:c1}).
Equalities (a)$=$(b) are one-line clock conversions
(\eqref{eq:kl}); the substantive claim is that the Fourier clock (c)
shows the same time, and the mechanism behind it is the localization of
the extremal frequency on the cycle-resonance line
$2^{k^*}\approx 3^m$ (Computational Fact~\ref{cf:numerics}), where the
large-deviation cost of tracking the line is exactly $I_0$.

\section{Relation to Tao's program and to the certificate barrier}
\label{sec:relation}

Tao's transport argument consumes the $3$-adic equidistribution of an
\emph{ensemble} of initial points; its pointwise failure is structural
(finite-window realization obstructions \cite{P1}). What an exponential
$\delta_m$ would buy is the \emph{effective constant} in that program:
the number of ternary levels needed per unit of density gain is
$1/I_0\approx 12.6$, and the boundary-layer analysis shows this constant
is sharp at the level of the Fourier input. Conversely, the barrier
theory of \cite{P1} prices the $2$-adic certificate side of the same
constant, and the extremal objects on both sides --- the all-ones tower
$-1 \bmod 2^m$ there, the pinned bridge on $2^{k}\approx 3^m$ here ---
are the two faces of the cycle-resonance line. We find it notable that
the obstruction isolated in Hypothesis~\ref{hyp:clocal} is again a
$\times 2\times 3$ statement: all roads from this problem's analytic
side appear to lead to the same rigidity frontier.

\section*{Data availability statement}
% TODO(author): confirm the repository URL / DOI before submission.
All code and data supporting this study --- the exponent-lattice
computations, the audits, the band-rate dynamic programs, the
certification drivers with their logs, and one self-contained script per
computational claim, with a claim-to-script map and SHA-256 manifest ---
are openly available in the reproduction repository
\texttt{collatz-barrier-repro}
(\url{https://github.com/hikaruwakaura/collatz-barrier-repro}); a
DOI-stamped archival snapshot of the exact version used here will be
deposited upon acceptance. The same material is included as the
supplementary bundle of this submission.

\appendix

\section{Data, scripts, and pre-registration log}\label{app:data}

\subsection*{Scripts} Each computational claim maps to one
self-contained exact-arithmetic script in the supplementary bundle:
the exponent-lattice recursion and full scans; the rescaled window
computation to $m=1200$; the first-passage decomposition audit
(identity to $10^{-14}$); the weighted-norm and capped-$\ell^1$ audits
($399$ levels, zero violations); the crossing/corner/pinned-bridge
decompositions; the $\ell^4$ collision functional; the all-$b$ product
DP for band rates; the cycle-law verification of
Theorem~\ref{thm:cstat}; and the certification driver of
Theorem~\ref{thm:instances} for $m=64,128,256$.

\subsection*{Pre-registration log (calibration)} Major numerical
discriminations were decided by criteria declared before the deciding
computation, and we report all outcomes, including negative ones, as a
calibration record. (1) An initial candidate for the decay constant,
$\rho=2^{3/2}/3$ (suggested by an apparent plateau at $m\le 60$), was
rejected by the pre-registered discrimination test at $m\le 200$: the
extremal exponent ratio $k^*/m$ crossed $3/2$ and continued toward
$\log_2 3$, and the local rate descended to the $I_0$ prediction; the
plateau was a finite-size crossing effect. (2) A pre-registered
main-term (positivity-based) lower-bound construction was rejected by
its own criterion (measured rate $I_0+0.0093$, below target), leading
to the crossing-term mechanism of Section~\ref{sec:mech}. (3) The
pre-registered two-sided band $[0.0790,0.0797]$ for the asymptotic rate
was replaced by the one-sided criterion (rate $\ge I_0-\epsilon$,
approaching $I_0$), on the structural ground that certified functionals
bound the rate from one side only; current data satisfy the one-sided
criterion. We keep this log because, in our experience, the two
rejections were the direct source of the mechanism results in
Section~\ref{sec:mech}.

\begin{thebibliography}{99}

\bibitem{Tao2022}
T.~Tao,
\emph{Almost all orbits of the Collatz map attain almost bounded values},
Forum Math. Pi \textbf{10} (2022), e12.

\bibitem{Terras1976}
R.~Terras,
\emph{A stopping time problem on the positive integers},
Acta Arith. \textbf{30} (1976), 241--252.

\bibitem{Lagarias1985}
J.~C. Lagarias,
\emph{The $3x+1$ problem and its generalizations},
Amer. Math. Monthly \textbf{92} (1985), 3--23.

\bibitem{Lagarias2010}
J.~C. Lagarias (ed.),
\emph{The Ultimate Challenge: The $3x+1$ Problem},
American Mathematical Society, Providence, RI, 2010.

\bibitem{KontorovichLagarias2010}
A.~V. Kontorovich and J.~C. Lagarias,
\emph{Stochastic models for the $3x+1$ and $5x+1$ problems and related
problems}, in: The Ultimate Challenge: The $3x+1$ Problem, Amer. Math.
Soc., 2010, 131--188.

\bibitem{Kesten1973}
H.~Kesten,
\emph{Random difference equations and renewal theory for products of
random matrices},
Acta Math. \textbf{131} (1973), 207--248.

\bibitem{Goldie1991}
C.~M. Goldie,
\emph{Implicit renewal theory and tails of solutions of random equations},
Ann. Appl. Probab. \textbf{1} (1991), 126--166.

\bibitem{Furstenberg1967}
H.~Furstenberg,
\emph{Disjointness in ergodic theory, minimal sets, and a problem in
Diophantine approximation},
Math. Systems Theory \textbf{1} (1967), 1--49.

\bibitem{BLMV2009}
J.~Bourgain, E.~Lindenstrauss, P.~Michel, and A.~Venkatesh,
\emph{Some effective results for $\times a\times b$},
Ergodic Theory Dynam. Systems \textbf{29} (2009), 1705--1722.

\bibitem{Erdos1979}
P.~Erd\H{o}s,
\emph{Some unconventional problems in number theory},
Math. Mag. \textbf{52} (1979), 67--70.

\bibitem{Lagarias2009}
J.~C. Lagarias,
\emph{Ternary expansions of powers of 2},
J. Lond. Math. Soc. (2) \textbf{79} (2009), 562--588.

\bibitem{Barina2021}
D.~Barina,
\emph{Convergence verification of the Collatz problem},
J. Supercomput. \textbf{77} (2021), 2681--2688.

\bibitem{Wirsching1998}
G.~J. Wirsching,
\emph{The Dynamical System Generated by the $3n+1$ Function},
Lecture Notes in Math. \textbf{1681}, Springer, Berlin, 1998.

\bibitem{P1}
H.~Wakaura,
\emph{Certificate complexity of uniform descent for the $3x+1$ problem},
preprint (2026).

\end{thebibliography}

\end{document}
